\documentclass[12pt,a4paper]{article}
\usepackage[margin=25mm]{geometry}
\usepackage{amsmath,amssymb,amsfonts}
\usepackage[most]{tcolorbox}
\usepackage{xcolor}
\usepackage{fancyhdr}
\usepackage{hyperref}
\definecolor{engineeringblue}{HTML}{1F4E79}
\definecolor{givenfill}{HTML}{EEF5FB}
\definecolor{solutionfill}{HTML}{FFF6D6}
\pagestyle{fancy}
\fancyhf{}
\lhead{EE 204 — Engineering mathematics}
\rhead{Homework 02}
\cfoot{\thepage}
\newcommand{\hwNumber}{02}
\newcommand{\hwTitle}{Linear systems and uncertainty}
\newtcolorbox{givenbox}{colback=givenfill,colframe=engineeringblue,title=Given,fonttitle=\bfseries}
\newtcolorbox{solutionbox}{colback=solutionfill,colframe=orange!70!black,title=Solution,fonttitle=\bfseries}

\begin{document}
\begin{titlepage}
\centering
\vspace*{18mm}
{\Large\bfseries Example University\\[3mm] School of Engineering}\par
\vfill
{\Huge\bfseries \hwTitle\\[5mm]}
{\Large Homework \hwNumber}\par
\vfill
{\large Student: \rule{65mm}{0.4pt}\\[4mm]
Instructor: Dr. A. Chen\\[4mm]
Due: 12 September 2026}\par
\vfill
\textit{Keep this page, the problem statements, and the worked reasoning in
one source file so the submission remains editable.}
\end{titlepage}

\section*{Instructions}
State assumptions, show units, and label the final answer.  The solution boxes
are included as worked examples; delete them or move them to a separate answer
key when submitting a student-only sheet.

\section*{Problem 1 — calibration line}
\begin{givenbox}
A load cell reads 2.4 mV at 0 N and 8.4 mV at 120 N.  Assume a linear response.
Find the sensitivity in mV/N and the expected reading at 75 N.
\end{givenbox}
\begin{solutionbox}
The slope is
\[
  s=\frac{8.4-2.4}{120-0}=0.05\;\mathrm{mV/N}.
\]
The model is $V=2.4+0.05F$, so $V(75)=6.15\;\mathrm{mV}$.
\end{solutionbox}

\section*{Problem 2 — a two-equation balance}
\begin{givenbox}
Solve $2x+y=11$ and $x-y=1$.  Interpret $x$ and $y$ as two component
contributions and report both values.
\end{givenbox}
\begin{solutionbox}
Adding the equations gives $3x=12$, hence $x=4$.  Substitution into the second
equation gives $y=3$.  The pair $(x,y)=(4,3)$ satisfies both balances.
\end{solutionbox}

\section*{Problem 3 — uncertainty}
\begin{givenbox}
The measured length is $L=25.0\pm0.2$ mm and the width is
$W=10.0\pm0.1$ mm.  Give a first-order estimate for the uncertainty in the
area $A=LW$.
\end{givenbox}
\begin{solutionbox}
For independent uncertainties, use
\[
  \Delta A \approx W\Delta L+L\Delta W
  =10.0(0.2)+25.0(0.1)=4.5\;\mathrm{mm^2}.
\]
Thus $A=250.0\pm4.5\;\mathrm{mm^2}$ under this conservative estimate.
\end{solutionbox}

\section*{Reflection}
Which assumption in one problem would you test first in a laboratory, and what
measurement would provide evidence for or against it?\\[12mm]
\rule{\linewidth}{0.4pt}\\[12mm]
\rule{\linewidth}{0.4pt}
\end{document}
